% ============================================================================
\clearpage
\section{De las oportunidades al producto}\label{sec:oportunidades}

La siguiente matriz cierra el recorrido completo: cada punto de fricción del journey map se
traza hasta la oportunidad de diseño que lo mitiga, la característica del producto que la
implementa, el escenario que la muestra en funcionamiento y la forma en que se verificará
cuando exista el prototipo. Es la continuación directa de la matriz de trazabilidad de la
actividad anterior, que conectaba pain points con características y personas.

\begin{uxtabla}{C{1.5cm}Z{1.05}Z{1.05}Z{0.9}C{2.2cm}}{Trazabilidad de cada punto de fricción hasta su verificación}
\uxheadrow \thd{Etapa} & \thd{Punto de fricción} & \thd{Oportunidad} & \thd{Característica del producto} & \thd{Cómo se verificará} \\ \midrule
2. Acceso & Permisos que no corresponden al rol & Espacio por rol y permiso mínimo suficiente & Acceso por permisos y roles & Tiempo de alta y de revocación \\
3. Búsqueda & Variables de nombre parecido & Fuente, propietario y vigencia en la tarjeta & Buscador unificado y catálogo & Aciertos en la tarea de localización \\
4. Validación & La definición vive en el correo de alguien & Ficha con definición, notas tribales y versiones & Catálogo de datos y trazabilidad & Tiempo hasta la validación \\
5. Profundización & Cifras sin fuente en la respuesta & Herramienta consultada visible y cita por cifra & Asistente con trazabilidad de la consulta & Confianza declarada en la prueba de usabilidad \\
6. Entrega & Copiar el dato pierde el contexto & Copiado con procedencia por omisión & Trazabilidad & Errores de atribución de la fuente \\
7. Reutilización & Nadie se entera del cambio de definición & Aviso dirigido a quien tiene la fuente guardada & Gobierno de metadatos y versiones & Uso de versiones retiradas \\
\end{uxtabla}

\subsection{Patrones de interacción que sostienen las oportunidades}\label{sec:patrones}

Las oportunidades de los cuatro mapas no son ideas sueltas: cada una se apoya en un patrón de
interacción que el equipo comprometió al definir el producto, y cada patrón tiene un
fundamento documentado. La tabla siguiente los reúne e indica en qué punto de qué recorrido se
manifiestan, de modo que ningún patrón quede como declaración de intenciones sin un momento
concreto donde el usuario lo perciba.

\begin{uxtabla}{L{3.4cm}Z{1.0}Z{1.0}}{Patrones de interacción comprometidos y dónde se manifiestan en los recorridos}
\uxheadrow \thd{Patrón} & \thd{Dónde se manifiesta} & \thd{Fundamento} \\ \midrule
Revelación progresiva & Etapas 3 y 5 del mapa de equipo; etapa 2 del mapa de Arturo Castañeda & Principio de diseño declarado en la actividad anterior: la complejidad aparece solo cuando la tarea la exige \\
Conocimiento tribal publicado junto al dato & Etapa 4 del mapa de equipo; escenario 4 completo & Agarwal et al. (2026): el conocimiento no documentado sobre las fuentes es lo que separa una consulta correcta de una equivocada \\
Visibilidad de la consulta de la inteligencia artificial & Etapa 5 del mapa de equipo; etapa 5 del mapa de Arturo Castañeda & Naik et al. (2026): en sistemas con agentes, mostrar qué herramienta se consultó es condición de confianza, no un adorno \\
Tarjetas con señal predictiva & Etapa 2 del mapa de Arturo Castañeda & La vista directiva debe entregar interpretación, no datos crudos \\
Espacios de trabajo por rol & Etapa 2 del mapa de equipo; etapa 1 del de Arturo Castañeda; etapa 3 del de Ximena Solís & El mismo producto debe presentar a cada perfil lo que su tarea requiere, incluidas las credenciales de integración \\
Estado compartido entre tablero y asistente & Etapa 5 del mapa de Arturo Castañeda & El resumen de la vista actual usa el estado del tablero como contexto de la consulta \\
Contrato de datos versionado con aviso de cambio & Etapas 4 y 6 del mapa de Ximena Solís; etapa 7 del mapa de equipo & Singh et al. (2026): el consumo gobernado exige contratos estables y comunicación anticipada de los cambios \\
\end{uxtabla}

Un patrón comprometido no aparece en ninguno de los cuatro recorridos: el \emph{streaming} con
detención real de la respuesta del asistente. Es deliberado y conviene declararlo: ninguno de
los seis escenarios se apoya en una conversación larga con la herramienta, de modo que la
cancelación no tiene todavía un momento en el que el usuario la perciba. Se retomará cuando la
actividad correspondiente aborde el diseño del asistente en alta fidelidad.

% ============================================================================
\section{Principios rectores y antiprincipios}\label{sec:principios}

El material de apoyo recomienda documentar los principios rectores del producto —una
descripción cualitativa de aquello que el producto defiende— y, en paralelo, los
antiprincipios, que declaran de forma explícita lo que el producto no pretende ser. Ambos
deben poder medirse, y su utilidad práctica es doble: alinean al equipo y ofrecen un criterio
neutral para rechazar funcionalidades que suenan atractivas pero desvían el producto.

Los siguientes se derivan directamente de los seis escenarios y del journey map, no de una
lluvia de ideas independiente.

\begin{uxtabla}{L{2.6cm}Z{1.15}Z{0.85}}{Principios rectores del portal y su forma de medición}
\uxheadrow \thd{Principio} & \thd{Definición} & \thd{Cómo se mide} \\ \midrule
Respaldado & Toda cifra que sale del portal puede citar su fuente, su fecha y su definición & Proporción de respuestas con procedencia completa \\
Simple primero & La consulta más frecuente se resuelve sin exponer filtros ni vocabulario técnico & Pasos hasta el resultado en la tarea de localización \\
Profundidad a un paso & Las opciones avanzadas están cerca y son descubribles por quien las necesita & Tasa de descubrimiento en la prueba de usabilidad \\
Sin sorpresas & Ningún cambio de definición o de esquema llega sin aviso previo ni versión anterior disponible & Avisos emitidos frente a cambios publicados \\
Explicable & El usuario puede ver qué consultó el sistema antes de decidir si cree en la respuesta & Visibilidad de la consulta en cada respuesta del asistente \\
\end{uxtabla}

\begin{uxtabla}{L{2.6cm}Z{1.15}Z{0.85}}{Antiprincipios: lo que el portal no pretende ser}
\uxheadrow \thd{Antiprincipio} & \thd{Definición} & \thd{Cómo se detecta} \\ \midrule
Caja negra & Respuestas sin fuente ni rastro de cómo se obtuvieron & Revisión de respuestas sin cita \\
Solo para expertos & Interfaz que supone conocimiento de la estructura de las bases & Tasa de éxito de los perfiles no técnicos \\
Autoritario & El sistema decide por el usuario cuando existe una ambigüedad legítima & Casos de desambiguación resueltos sin consultarlo \\
Silencioso & Cambios que rompen procesos aguas abajo sin notificación previa & Incidentes originados por cambios no anunciados \\
\end{uxtabla}

% ============================================================================
\clearpage
\section{Referencias}

\begin{uxreferencias}
Agarwal, S., Biswal, A., Zeighami, S., Cheung, A., Gonzalez, J., y Parameswaran, A. G. (2026).
Arming data agents with tribal knowledge. \textit{arXiv}. https://arxiv.org/abs/2602.13521

Baxter, K., Courage, C., y Caine, K. (2015). \textit{Understanding your users: A practical
guide to user research methods} (2.\textsuperscript{a} ed.). Morgan Kaufmann.

Hartson, R., y Pyla, P. S. (2012). \textit{The UX book: Process and guidelines for ensuring a
quality user experience}. Morgan Kaufmann.

International Organization for Standardization. (2018). \textit{Ergonomics of human-system
interaction — Part 11: Usability: Definitions and concepts} (ISO 9241-11:2018).

McInerney, P. (2003). Plantilla de escenarios [citado en Baxter, K., Courage, C., y Caine, K.
(2015), \textit{Understanding your users: A practical guide to user research methods},
2.\textsuperscript{a} ed., Morgan Kaufmann].

Naik, S., Passi, S., Vorvoreanu, M., Saponas, S., y Hall, A. (2026). ``So there's a Catch-22
here'': How early adopters who build multi-agent LLM systems conceptualize transparency.
\textit{arXiv}. https://arxiv.org/abs/2606.08323

Osterwalder, A., y Pigneur, Y. (2010). \textit{Business model generation: A handbook for
visionaries, game changers, and challengers}. John Wiley \& Sons.

Ramírez Mejía, A. I. (2023). \textit{Grocery shopping app: A guided UX case study. Part 2.
Scenarios and journey map} [Diapositivas]. TC4032 Experiencia del usuario y diseño de
interfaces, Maestría en Inteligencia Artificial Aplicada, Instituto Tecnológico y de Estudios
Superiores de Monterrey.

Singh, G., Kavehzadeh, P., Xia, J., Fu, X.-Y., Bouvier Tremblay, J., Laskar, M. T. R., Lum, V.,
y Bhushan TN, S. (2026). Beyond text-to-SQL: An agentic LLM system for governed enterprise
analytics APIs. \textit{arXiv}. https://arxiv.org/abs/2605.21027

Walter, S. (2022). \textit{User journey mapping: Visualize user research, brainstorm
opportunities, and solve problems}. SitePoint.
\end{uxreferencias}

